\chapter{Sumber Cahaya dan Perambatannya}\label{ch:g7:light-sources-propagation}

Mengapa kamu dapat membaca halaman ini? Sebuah lampu bersinar ---
tetapi halamannya bukan lampu, namun matamu menerima \omterm{def:g1:light-and-shadows:source}{cahaya} dari
setiap baris cetakannya. Di antara benda yang membuat \omterm{def:g1:light-and-shadows:source}{cahaya} dan benda
yang sekadar meneruskannya terletak perbedaan yang tidak dapat
ditinggalkan optika; dan bersamanya datang jawaban penuh bagi
pertanyaan yang tertua di dalam buku ini: apa yang diperlukan untuk
melihat sesuatu?

\section{Pembuat dan penerus}

\begin{definition}[Sumber primer]\label{def:g7:light-sources-propagation:primary}
Sebuah \emph{sumber primer}\index{sumber primer} membuat \omterm{def:g1:light-and-shadows:source}{cahayanya}
sendiri: Matahari dan \omterm{def:g4:solar-system:star}{bintang}, nyala api, filamen yang berpijar, baja
yang berpijar putih di tempat penempaan, kunang-kunang, kilat.
\omterm{def:g1:light-and-shadows:source}{Cahayanya} ada, entah ada entah tidak ada yang lain yang bersinar ---
padamkan dunianya, dan sumber primer tetap berpijar.
\end{definition}

\begin{definition}[Benda pembaur]\label{def:g7:light-sources-propagation:diffusion}
Setiap benda terlihat yang lain adalah \emph{benda
pembaur}\index{pembauran}: ia menerima \omterm{def:g1:light-and-shadows:source}{cahaya} lalu mengirim ulang ---
\emph{membaurkan} --- sebagian darinya ke segala arah. Halamannya,
apelnya, \omterm{def:g2:moon-in-the-sky:moon}{Bulan}, wajahmu sendiri: tidak satu pun membuat kerlip
miliknya sendiri; masing-masing menghamburkan \omterm{def:g1:light-and-shadows:source}{cahaya} yang diterimanya
dengan murah hati ke segala sisi. Pembaur di dalam gelap sekadar tidak
terlihat.
\end{definition}

\begin{omfigure}
\begin{tikzpicture}[scale=0.8]
  % lampunya, sebuah sumber primer
  \fill[omThm] (0.7,2.7) circle (0.3);
  \foreach \a in {0,45,...,315} {
    \draw[omThm, thick] (0.7,2.7) ++(\a:0.42) -- ++(\a:0.2);
  }
  \node[above, font=\small] at (1.1,3.35) {lampu: sumber primer};
  % cahayanya berjalan menuju apelnya
  \draw[-Stealth, omThm, thick] (1.35,2.75) -- (3.55,2.15);
  \draw[-Stealth, omThm, thick] (1.35,2.45) -- (3.45,1.65);
  \draw[-Stealth, omThm, thick] (1.25,2.15) -- (3.35,1.15);
  % apelnya menghamburkannya ke segala sisi
  \draw[thick, fill=omMeth, fill opacity=0.5] (5.2,1.2) circle (0.5);
  \foreach \a in {15,60,105,150,195,240,285,330} {
    \draw[-Stealth, omDef, thick] (5.2,1.2) ++({\a}:0.6)
      -- ++({\a}:1.05);
  }
  \node[below, font=\small] at (5.2,-0.5) {apel: pembaur};
  \node[right, font=\small, omDef] at (7.35,2.15)
    {cahaya yang dibaurkan pergi ke segala sisi:};
  \node[right, font=\small, omDef] at (7.35,1.55)
    {setiap mata di ruangan itu mendapat sebagian};
\end{tikzpicture}

{\small Pembagian kerjanya: lampunya membuat \omterm{def:g1:light-and-shadows:source}{cahaya}; apelnya, yang
menerimanya, menghamburkan sebagian ke segala arah --- itulah sebabnya
setiap orang di sekeliling mejanya melihat apel yang sama.}
\end{omfigure}

\begin{example}[Pembauran bukan pencerminan]\label{ex:g7:light-sources-propagation:notmirror}
Baik \omterm{def:g5:mirrors-reflection:reflection}{cermin} maupun halaman mengirim kembali \omterm{def:g1:light-and-shadows:source}{cahaya} yang diterimanya
--- tetapi dengan tata cara yang berlawanan. \omterm{def:g5:mirrors-reflection:reflection}{Cerminnya} memantulkan ke
\emph{satu} arah yang tertib (sudut yang sama --- aturan pantulanmu
yang lama), sehingga gambarnya terjaga; kekasaran mikroskopis
halamannya menyemprotkan \omterm{def:g1:light-and-shadows:source}{cahaya} ke \emph{mana-mana}, sehingga
merusakkan gambarnya tetapi menawarkan dirinya kepada setiap mata
sekaligus. Itulah sebabnya \omterm{def:g5:mirrors-reflection:reflection}{cermin} memperlihatkan lampunya, sedangkan
halamannya sekadar memperlihatkan dirinya, dari mana pun kamu berdiri.
\end{example}

\begin{example}[Berkas Bulan, ditutup]\label{ex:g7:light-sources-propagation:moon}
Putusan masa kecilnya kini mempunyai kata resminya: \omterm{def:g2:moon-in-the-sky:moon}{Bulan} bukan \omterm{def:g7:light-sources-propagation:primary}{sumber
primer}; ia \omterm{def:g7:light-sources-propagation:diffusion}{benda pembaur} seukuran \omterm{def:g4:solar-system:planet}{planet} --- batu kelabu yang
menghamburkan \omterm{def:g1:light-and-shadows:source}{cahaya} matahari ke segala arah, termasuk ke arah kita.
Jadi \omterm{def:g1:light-and-shadows:source}{cahaya} bulan adalah sinar matahari bekas, dan ``sisi gelapnya''
sekadar separuh yang berkasnya, pada saat itu, tidak memuat \omterm{def:g1:light-and-shadows:source}{cahaya}
untuk diteruskan.
\end{example}

\section{Melihat, dituntaskan}

\begin{proposition}[Kedua syarat penglihatan]\label{prop:g7:light-sources-propagation:vision}
Sebuah mata melihat sebuah benda tepat ketika
\begin{enumerate}
  \item \omterm{def:g1:light-and-shadows:source}{cahaya} \emph{meninggalkan} bendanya --- dibuat, kalau
        bendanya \omterm{def:g7:light-sources-propagation:primary}{sumber primer}; dibaurkan, kalau ia penerus --- dan
  \item \omterm{def:g1:light-and-shadows:source}{cahaya} itu berjalan menurut garis lurus yang tak terhalang
        masuk ke matanya.
\end{enumerate}
Patahkan salah satu syaratnya --- pembaur yang tidak tersinari, atau
tembok yang melintang pada garis pandangnya --- dan bendanya tidak
terlihat. Setiap teka-teki keterlihatan di dalam buku ini terurai
menjadi kedua pemeriksaan ini.
\end{proposition}

\begin{example}[Berkas yang dibuat terlihat]\label{ex:g7:light-sources-propagation:beam}
Di \omterm{def:g2:air-around-us:air}{udara} yang bersih, \omterm{def:g6:light-propagation:ray}{berkas} senter yang menyeberangi ruangan tidak
dapat dilihat dari samping --- \omterm{def:g2:air-around-us:air}{udara} meneruskan hampir tidak ada, dan
tidak ada \omterm{def:g1:light-and-shadows:source}{cahaya} dari berkasnya yang berbelok ke arah matamu.
Kibaskan kain yang berdebu: berkasnya melompat ke dalam pandangan.
Tiap bintik debu adalah pembaur mungil, yang menangkap \omterm{def:g1:light-and-shadows:source}{cahaya} berkasnya
lalu menghamburkan sebagian ke arahmu --- ribuan bintik menggambarkan
garis lurus berkasnya untukmu. Kabut, embun, dan asap menjalankan jasa
yang sama bagi \omterm{def:g6:light-propagation:ray}{berkas} mercusuar dan proyektor bioskop.
\end{example}

\section{Tiga mantel bagi cahaya}

\begin{definition}[Bening, baur, buram]\label{def:g7:light-sources-propagation:coats}
Bahan bertemu \omterm{def:g1:light-and-shadows:source}{cahaya} dengan tiga tata cara:
\begin{enumerate}
  \item \emph{bening}\index{bening}: \omterm{def:g1:light-and-shadows:source}{cahaya} menyeberang hampir tanpa
        terganggu, sambil menjaga ketertibannya --- kaca jernih, air
        yang tenang, \omterm{def:g2:air-around-us:air}{udara}; gambarnya lewat dengan utuh;
  \item \emph{baur}\index{baur}: \omterm{def:g1:light-and-shadows:source}{cahaya} menyeberang, tetapi teracak
        --- kaca beku, kertas kalkir, tirai yang tipis; terangnya
        lewat, gambarnya tidak;
  \item \emph{buram}\index{buram}: \omterm{def:g1:light-and-shadows:source}{cahaya} dihentikan --- diserap atau
        dibaurkan kembali --- kayu, batu, logam, tanganmu; di balik
        benda yang buram terletak \omterm{def:g1:light-and-shadows:shadow}{bayang-bayang}.
\end{enumerate}
\end{definition}

\begin{method}[Memilah bahan dengan cahaya senter]\label{met:g7:light-sources-propagation:sort}
Sebuah ruangan gelap, sebuah senter, sehalaman cetakan, dan
contohnya di antara keduanya:
\begin{enumerate}
  \item peganglah contohnya di atas cetakannya dengan senter di
        belakangnya;
  \item cetakannya terbaca menembusnya? --- \omterm{def:g7:light-sources-propagation:coats}{bening};
  \item tampak nyala tetapi tidak tampak hurufnya? --- \omterm{def:g7:light-sources-propagation:coats}{baur};
  \item tampak kegelapan, dan \omterm{def:g1:light-and-shadows:shadow}{bayang-bayang} di seberangnya? --- \omterm{def:g7:light-sources-propagation:coats}{buram}.
\end{enumerate}
Ujilah dengan jujur: banyak bahan berpindah regu menurut tebalnya ---
selapis air itu \omterm{def:g7:light-sources-propagation:coats}{bening}, laut dalam segelap malam; selembar kertas itu
\omterm{def:g7:light-sources-propagation:coats}{baur}, buku yang tertutup \omterm{def:g7:light-sources-propagation:coats}{buram}.
\end{method}

\begin{omfigure}
\begin{tikzpicture}[scale=0.8]
  \foreach \x/\name in {0.5/bening, 4.4/baur,
      8.3/buram} {
    \fill[omThm] (\x,1.5) circle (0.22);
    \draw[thick] (\x+1.3,0.4) rectangle (\x+1.7,2.6);
    \node[below, font=\small] at (\x+1.6,-0.0) {\name};
  }
  \draw[-Stealth, omThm, thick] (0.75,1.5) -- (1.78,1.5);
  \draw[-Stealth, omThm, thick] (2.22,1.5) -- (3.4,1.5);
  \draw[-Stealth, omThm, thick] (4.65,1.5) -- (5.68,1.5);
  \foreach \a in {28,10,-12,-30} {
    \draw[-Stealth, omThm, thin] (6.12,1.5) -- ++({\a}:1.1);
  }
  \draw[-Stealth, omThm, thick] (8.55,1.5) -- (9.58,1.5);
  \draw[line width=3pt, omExo] (10.02,0.55) -- (10.02,2.45);
  \node[right, font=\footnotesize, omExo] at (10.1,0.8) {sisi bayang-bayang};
\end{tikzpicture}

{\small Tiga nasib bagi sebuah sinar: menembus dengan tertib;
menembus tetapi teracak; dihentikan, dengan \omterm{def:g1:light-and-shadows:shadow}{bayang-bayang} di belakang.}
\end{omfigure}

\begin{example}[Ketiga regu milik arsitektur]\label{ex:g7:light-sources-propagation:architecture}
Pembangun mengerahkan ketiganya dengan sengaja. Kaca jendela yang
\omterm{def:g7:light-sources-propagation:coats}{bening}: pemandangan dan \omterm{def:g1:light-and-shadows:source}{cahaya} sekaligus. Kaca kamar mandi yang \omterm{def:g7:light-sources-propagation:coats}{baur}
dan sekat kantor yang beku: \omterm{def:g1:light-and-shadows:source}{cahaya} siang masuk, kesendirian terjaga
--- terang tanpa gambar tepat menjadi uraian tugasnya. Tembok dan
daun jendela yang \omterm{def:g7:light-sources-propagation:coats}{buram}: kegelapan sesuai permintaan. Kap lampu adalah
kebauran yang mengabdi pada kenyamanan: silau filamen yang berkobar
teracak menjadi nyala yang lembut dan tanpa sumber.
\end{example}

\begin{remark}[Tidak ada yang diteruskan, tidak ada yang terlihat]\label{rem:g7:light-sources-propagation:space}
Teka-teki: seorang antariksawan di antara \omterm{def:g4:solar-system:star}{bintang} memandang \emph{ke
samping} menyeberangi kehampaan yang dibanjiri sinar matahari. Apa
yang dilihatnya? Kehitaman --- padahal \omterm{def:g1:light-and-shadows:source}{cahaya} matahari membanjir lewat
di depan matanya. Antariksa tidak membawa debu yang patut disebut:
tidak ada yang membaurkan \omterm{def:g1:light-and-shadows:source}{cahaya} yang lewat itu ke arahnya, dan \omterm{def:g1:light-and-shadows:source}{cahaya}
yang tidak masuk ke mata tidak melukis apa-apa. Langit siang berwarna
biru justru karena \omterm{def:g2:air-around-us:air}{udara} kita \emph{memang} meneruskan \omterm{def:g1:light-and-shadows:source}{cahaya} matahari
turun kepada kita --- \omterm{def:g7:light-sources-propagation:diffusion}{pembauran} lembut sebanyak seluruh \omterm{def:g4:solar-system:planet}{planet} di atas
kepala. Tanpa \omterm{def:g2:air-around-us:air}{udara}, langit hitam, bahkan pada tengah hari: para
antariksawan \omterm{def:g2:moon-in-the-sky:moon}{Bulan} berdiri di dalam sinar matahari di bawah kubah
hitam tanpa \omterm{def:g4:solar-system:star}{bintang}.
\end{remark}

\section{Latihan}

\begin{exercise}[$\star$]\label{exo:g7:light-sources-propagation:1}
Pilahlah: nyala lilin; \ \omterm{def:g2:moon-in-the-sky:moon}{Bulan}; \ seekor kunang-kunang; \ halaman ini;
\ lampu lalu lintas yang merah; \ \omterm{def:g4:solar-system:planet}{planet} Venus pada waktu senja.
\end{exercise}

\begin{exercise}[$\star$]\label{exo:g7:light-sources-propagation:2}
Apa yang dilakukan \omterm{def:g7:light-sources-propagation:diffusion}{benda pembaur} terhadap \omterm{def:g1:light-and-shadows:source}{cahaya} yang diterimanya ---
dan apa jadinya ia di dalam ruangan yang gelap?
\end{exercise}

\begin{exercise}[$\star$]\label{exo:g7:light-sources-propagation:3}
Nyatakan kedua syarat penglihatannya, lalu pakailah: mengapa kamu
tidak dapat melihat permukaan \emph{milik} \omterm{def:g5:mirrors-reflection:reflection}{cermin} yang sempurna
bersih, melainkan hanya benda yang dipantulkannya? (Hati-hati ---
pertanyaan yang halus; jawablah dengan sederhana: apa yang ditolak
kilapnya untuk dilakukan terhadap \omterm{def:g1:light-and-shadows:source}{cahaya}?)
\end{exercise}

\begin{exercise}[$\star$]\label{exo:g7:light-sources-propagation:4}
Mengapa setiap orang di sekeliling mejanya melihat apel yang sama,
sementara hanya satu kedudukan mata yang beruntung yang menangkap
pantulan lampunya pada sebuah sendok?
\end{exercise}

\begin{exercise}[$\star$]\label{exo:g7:light-sources-propagation:5}
Mengapa \omterm{def:g6:light-propagation:ray}{berkas} senter tidak terlihat dari samping di \omterm{def:g2:air-around-us:air}{udara} yang
bersih, dan terlihat di dalam kabut? Sebutkan prosesnya.
\end{exercise}

\begin{exercise}[$\star$]\label{exo:g7:light-sources-propagation:6}
Golongkanlah: kaca jendela; \ kertas kalkir; \ sebuah batu bata; \ air
jernih yang dangkal; \ selubung sebuah \omterm{def:g3:first-electric-circuit:bulb}{bohlam} beku; \ kertas timah
aluminium.
\end{exercise}

\begin{exercise}[$\star$]\label{exo:g7:light-sources-propagation:7}
Tugas apa yang dijalankan kebauran pada jendela kamar mandi dan pada
kap lampu? Apa yang lewat, apa yang dihentikan?
\end{exercise}

\begin{exercise}[$\star$]\label{exo:g7:light-sources-propagation:8}
\omterm{def:g2:moon-in-the-sky:moon}{Bulan} tersinari matahari --- namun para antariksawan bulan melihat
langit hitam pada tengah hari. Jelaskan dengan
\cref{rem:g7:light-sources-propagation:space}, lalu katakan mengapa
langit tengah hari Bumi justru biru cerah.
\end{exercise}

\begin{exercise}[$\star\star$]\label{exo:g7:light-sources-propagation:9}
Selembar kertas tunggal bernyala hangat di depan sebuah senter; buku
yang tertutup menghentikan \omterm{def:g1:light-and-shadows:source}{cahayanya} seluruhnya. Apa yang dilakukan
tebalnya terhadap keanggotaan regu sebuah bahan? Berikan contoh kedua
tentang bahan yang berpindah regu.
\end{exercise}

\begin{exercise}[$\star\star$]\label{exo:g7:light-sources-propagation:10}
Muslihat bioskop: \omterm{def:g6:light-propagation:ray}{berkas} proyektor tidak terlihat di dalam aula yang
bersih, namun film memperlihatkannya mengiris \omterm{def:g2:air-around-us:air}{udara} dengan dramatis.
Apa yang harus ditambahkan kru panggung ke \omterm{def:g2:air-around-us:air}{udaranya}, dan fisika apa
yang sedang mereka beli?
\end{exercise}

\begin{exercise}[$\star\star$]\label{exo:g7:light-sources-propagation:11}
Kamu melihat wajah seorang kawan di dekat api unggun. Telusurilah
perjalanan penuh \omterm{def:g1:light-and-shadows:source}{cahayanya} --- sumber, \omterm{def:g7:light-sources-propagation:diffusion}{pembauran}, garis lurus --- lalu
hitunglah berapa banyak pembaur yang berdiri pada rantainya kalau kamu
lalu \emph{juga} melihat wajah itu terpantul di danaunya.
\end{exercise}

\begin{exercise}[$\star\star\star$]\label{exo:g7:light-sources-propagation:12}
Pertunjukan \omterm{def:g1:light-and-shadows:source}{cahaya} milik alam sendiri: pada hari yang cerah, langit
yang jauh dari Matahari berwarna biru, cakram Mataharinya putih
menyilaukan --- dan pada waktu \omterm{def:g1:day-and-night:daynight}{matahari terbenam} Matahari yang rendah
berubah jingga sementara awan di atasnya menangkap warna merah muda.
Dengan memakai hanya ``\omterm{def:g2:air-around-us:air}{udara} membaurkan sedikit \omterm{def:g1:light-and-shadows:source}{cahaya} yang lewat
menembusnya, dan makin banyak makin panjang jalannya'', jelaskan
mengapa matahari yang rendah dirampok sebagian \omterm{def:g1:light-and-shadows:source}{cahayanya} --- dan siapa
yang menerima jatah curiannya. (Kisah warnanya yang penuh adalah bab
spektrum tahun depan; bernalarlah di sini hanya dengan jalan lebih
panjang, \omterm{def:g7:light-sources-propagation:diffusion}{pembauran} lebih banyak.)
\end{exercise}

\section{Soal: Menerangi Museum}

\begin{problem}[{Soal akhir pekan --- aula pahatan yang baru; giliran
malam seorang perancang pencahayaan; sumber, pembaur, dan tiga jenis
kaca}]\label{pb:g7:light-sources-propagation:1}
Aula baru museumnya dibuka hari Senin: sebuah patung pualam,
lukisan di dindingnya, sebuah lemari pajang kaca, dan seorang kurator
yang gugup. Kamu membantu perancang pencahayaannya.

\medskip\noindent\textbf{Bagian I --- Asas pertama di tempatnya.}
\begin{enumerate}
  \item Kuratornya bertanya mengapa patungnya, yang ``begitu putih
        cemerlang'', masih memerlukan lampu sama sekali. Jawablah
        dengan kedua syarat penglihatannya.
  \item Sebutkan \omterm{def:g7:light-sources-propagation:primary}{sumber primer} dan pembaur di dalam aulanya begitu
        lampunya menyala: lampu, patung, lukisan, pengunjung, dan kaca
        lemarinya.
  \item Pualamnya seakan-akan bernyala dari mana-mana, tanpa bintik
        yang menyilaukan. Tata cara pengembalian \omterm{def:g1:light-and-shadows:source}{cahaya} yang mana yang
        sedang bekerja, dan mengapa setiap kedudukan pengunjung
        menerimanya?
  \item Sebaliknya, sekeping papan kuningan yang berkilap melemparkan
        satu pantulan lampu yang menyilaukan. Sebutkan tata cara yang
        satu lagi, lalu nyatakan aturan yang mengatur ke mana silau
        itu pergi.
\end{enumerate}

\medskip\noindent\textbf{Bagian II --- Menjinakkan silaunya.}
\begin{enumerate}[resume]
  \item \omterm{def:g3:first-electric-circuit:bulb}{Bohlam} telanjang menyilaukan mata pengunjung. Perancangnya
        menyarungkan tiap \omterm{def:g3:first-electric-circuit:bulb}{bohlam} ke dalam bola kaca beku. Regu bahan
        yang mana yang ditempati bolanya, dan apa yang dilakukannya
        terhadap \omterm{def:g1:light-and-shadows:source}{cahaya} filamennya?
  \item Lukisannya harus tersinari, tetapi tidak boleh ada lampu yang
        tampak \emph{terpantul} pada peliturnya dari tempat menonton
        yang utama. Dengan memakai aturan \omterm{def:g5:mirrors-reflection:reflection}{cerminnya}, di mana
        perancangnya menempatkan lampunya terhadap lukisan dan
        penontonnya?
  \item Kaca jernih lemari pajangnya hampir tidak terlihat --- sampai
        lengan baju putih seorang pengunjung muncul di dalamnya
        seperti hantu. Mengapa kaca jernih sekaligus meneruskan dan
        samar-samar mencerminkan? Apa yang dicapai ceruk berdinding
        hitam di belakang lemarinya, yang dibuat perancangnya?
  \item Bintik debu kadang-kadang hanyut menembus \omterm{def:g6:light-propagation:ray}{berkas} lampu
        sorotnya, dan menggambar benang yang terang di \omterm{def:g2:air-around-us:air}{udara}.
        Kuratornya: ``Indah! Pertahankan.'' Perancangnya: ``Itu
        berarti \omterm{def:g2:air-around-us:air}{udaranya} kotor.'' Siapa yang benar tentang fisikanya,
        dan benangnya itu apa?
\end{enumerate}

\medskip\noindent\textbf{Bagian III --- Uji malamnya.}
\begin{enumerate}[resume]
  \item Lampu padam, daun jendela tertutup rapat: penjaganya
        melaporkan aulanya ``hitam sama sekali --- bahkan patung
        putihnya pun tidak''. Jelaskan dengan kedua syaratnya mengapa
        keputihan tidak berdaya di dalam gelap.
  \item Diusulkan sebuah jendela atap: kaca beku di atas kepala. Apa
        yang akan diantarkannya pada siang hari, dan apa yang akan
        ditolaknya untuk diantarkan (pemandangan awan, silau matahari
        langsung)?
  \item Kuratornya takut \omterm{def:g1:light-and-shadows:source}{cahaya} matahari memudarkan lukisannya lalu
        meminta ``\omterm{def:g1:light-and-shadows:source}{cahaya} tanpa matahari sama sekali''. Perancangnya
        tersenyum: pada siang hari, bahkan \omterm{def:g1:light-and-shadows:source}{cahaya} lembut jendela yang
        menghadap ke utara pun \emph{adalah} sinar matahari. Jelaskan
        --- apa yang telah meneruskannya?
  \item Tulislah laporan perancangnya: tiga kalimat --- satu tentang
        mengapa \omterm{def:g7:light-sources-propagation:diffusion}{pembauran} adalah sahabat terbaik museumnya, satu
        tentang di mana pantulan bertata cara \omterm{def:g5:mirrors-reflection:reflection}{cermin} harus diawasi,
        dan satu tentang regu kaca yang dipilih bagi tiap tugasnya.
\end{enumerate}
\end{problem}
